% !TeX root = ../../Book3.tex
\chapter*{序言}
\addcontentsline{toc}{chapter}{序言}

复变量的微分看起来只比实变量多了一种方向，结论却十分不同。一个全纯函数在小圆上的值决定圆内的值，局部的零点与极点能够控制闭曲线积分，区域的形状又会改变函数的边界行为。沿着这些联系继续走，便会遇到覆盖空间、曲率、椭圆函数和模形式，最终回到关于整数和素数的问题。

\ProseParagraphBreak

本卷希望把这条路线写成可以逐步学习的教材。第一编从幂级数与 Cauchy 理论开始，给出基础复分析的系统复习；已经学过一学期复分析的读者，可先检查章首目标和习题，再决定哪些证明需要重读。其余部分依次讨论共形几何、势论、整函数、椭圆与模函数，以及 Dirichlet 级数和素数定理。这里需要多元微积分、线性代数、基本测度积分和度量空间知识；抽象拓扑、覆盖与曲面语言则在使用前建立或在附录补充。

\ProseParagraphBreak

第二编的起点是 Riemann 映照定理。接着，Schwarz--Pick 估计把映射变成距离比较，边界对应和极值长度显示区域形状的限制，解析延拓与覆盖又把局部函数放进整体空间。统一化说明共形类型之间的差别，负曲率方法随后通向 Picard 定理。读这一编时，尤其要把区域、覆盖空间与映射的定义域分别画清；同一个局部表达式放在不同空间中，可能代表不同的全局对象。

\ProseParagraphBreak

第三编把全纯函数的实部和模的对数作为主要对象。调和函数、次调和函数、Perron 方法与 Green 函数提供边值和势论工具；Poisson--Jensen 公式再把边界平均与零点联系起来。第四编据此讨论整函数增长、乘积和分解，并以 Gamma 与 Mellin 方法展示解析延拓和渐近计算。这里的收敛域、无穷乘积正规收敛以及对数分支，都须与形式恒等式同时核验。

\ProseParagraphBreak

第五编从格上的双周期函数走向模形式与 theta 函数。前面的留数、紧性和零点计数在新的对称性下继续发挥作用，Fourier 系数则逐渐显出算术含义。第六编研究算术函数、Dirichlet 级数与 zeta 函数，最后证明素数定理。素数计数不是一次留数计算的直接产物：必须先得到适当的解析延拓与边界零点排除，再证明把解析信息转为平均渐近的步骤。

\ProseParagraphBreak

本卷可以按不同目的分段阅读。共形几何路线以第一、二编为主，并查阅 Riemann 曲面与拟共形附录；势论和整函数路线接读第三、四编；解析数论路线则在掌握留数、正规收敛和整函数工具后，进入第五、六编。各编导读给出更具体的依赖图。附录中的 Riemann--Roch、可测 Riemann 映射等深层理论有明确的陈述与证明范围，不能因它们出现在附录就忽略其假设，也不必在首次阅读时把所有分支同时学完。

\ProseParagraphBreak

书中统一用保向的意义理解共形映射，区分单连通、平面 Green 域与普适覆盖的共形类型。符号相近的对象会在定义时分别说明，例如上半平面与圆盘的双曲度量、不同尺度的 theta 变量、已完成的 zeta 函数及其未完成形式。与第二卷联系的 Fourier 计算也保留 \(2\pi\) 约定，避免在变量替换中丢失常数。

\ProseParagraphBreak

每章的例子与习题不仅用于熟悉计算，也用于检查结论的边界。一个穿孔域、一个多值对数或一列在边界聚集的零点，常能说明为何某个定理必须要求连通性、单值性或内部紧性。解答供核对核心方法，不宜先以结论是否正确判断整份证明。正文中的中心论证已在相应位置展开；采用证明概要的深层结果，另标来源和省略内容。

\ProseParagraphBreak

基础复分析可伴读 Conway 的 \cite{Conway1978Functions} 或 Stein 与 Shakarchi 的 \cite{Stein2003Complex}；共形不变量、势论与模函数可分别从 \cite{Ahlfors2010Conformal,Ransford1995Potential,Apostol1990Modular} 继续。后续专著往往压缩了本书展开的基础步骤，也常采用更一般的空间与更强的结论。比较时先核对所讨论的对象和假设，再比较公式与证明，会少走许多弯路。

\ProseParagraphBreak

学习复分析时，可以反复追问一个问题：眼前的局部信息是怎样成为整体限制的？有时答案来自积分，有时来自拓扑，有时来自增长或曲率。把那一步真正写清，便更容易理解为何看似相距很远的共形映射与素数计数，会在同一卷中相遇。

\bigskip
\begin{flushright}
R. EverStarine

2026年10月11日
\end{flushright}
